\section{My impact}

\textbf{3a. How my contributions drove the team forward.} \rtag{M16.a}\rtag{M17.a} The test is not what I built but whether, in its absence, the team could have reached submission. I think the honest answer is no on several items; the table below lists the ones where I can point to a measurable outcome.

{\footnotesize
\begin{tabularx}{\textwidth}{@{}p{3.4cm} X p{4.3cm}@{}}
\toprule
\textbf{Gap identified} & \textbf{Action I took (dated, past tense)} & \textbf{Measurable outcome} \\
\midrule
No flight data after wk16 cancelled day 1 & Built SITL simulation framework (\texttt{simple\_simulator.py}, 2508 lines) with GPS drift, camera shake, CV throttle, TFLite backend, god-view, keyboard flight & Replaced $\sim$50 physical flights; the only source of verifiable mission data the team had \cite{complaintletter} \\
No progressive path from bench $\to$ flight & Built 5-step numbered test ladder (\texttt{tests/flight/0a\_cube\_commands} $\to$ \texttt{4\_detect\_and\_center}) + 6 gap-fill scripts in one session & 41 categorised test scripts, documented protocol in \texttt{FLIGHT\_DAY\_TESTS.md}; pilot could follow it without me \\
\texttt{main.py} unreadable to teammates (wk17 pushback, 11-state FSM) & Slept on feedback; next day refactored 1411$\to$754 lines, extracted 6 modules, added 146 \texttt{pytest} tests \cite{mainrefactor} & Teammates able to read state machine; re-asserted wk17 as the start of teachable architecture \\
Detection maths wrong on tilted camera & Iterated ray-tracing tilt compensation through 12 commits [\texttt{ad207ae}$\to$\texttt{831c47d}], verified against a 1440-test sweep using standard pinhole projection maths & GPS estimation error progression: 7~m $\to$ 3.5~m $\to$ 1.8~m $\to$ 1.5~m CEP50 \cite{gpsestaccuracy} \\
Robin's script needed camera without duplicating CV code & Built \texttt{passive\_watch\_2.py} with \texttt{-{}-robin} flag, \texttt{cv\_mode} polling, 2-image output, JSON schema README \cite{robinhandoff} [\texttt{a5b1ab7, 885cee4, 36b6069}] & Robin integrated in one session without reading my CV code; two modules run in parallel without conflict \\
Demetro's FDR slides not ready wk21 & Built his 2 full-screen HTML slides + wrote 465-word speaking scripts, iterated 4$\times$ [\texttt{32fb0c5, b2a9cd4, ba667f2, 9b826f8, 7920e5a, 794ef5c}] & Demetro delivered 3.1-minute talk from polished scripts; his time went to the aerodynamics section instead \\
Hardware failures unacknowledged by course & Drafted team complaint letter, 6 iterations over 4 days, measured tone (``we are not assigning blame, we are documenting'') [\texttt{26e3255}$\to$\texttt{4b0fd91}] & Letter submitted as team advocacy on behalf of all 5; assessment context formally on record \\
R01--R12 compliance gaps in wk22 audit & Walked brief one-by-one, closed 12 gaps in one session (runtime geofence \texttt{pointPolygonTest} $\to$ RTL, 50~m altitude cap, 5~m take-off landing distance check, landing distance print, detection photo logging, MIT LICENSE file) [commit \texttt{ce5c036}] & All 12 requirements MET in simulation; report can claim verifiable compliance \\
\bottomrule
\end{tabularx}}

\textit{What I would not want the reader to mistake.} \rtag{M16.d-team} The simulation, test infrastructure, and integration discipline drove the team forward in the sense that they kept the project alive. They did not make up for the absence of real flight data. The gap between ``the simulator says X'' and ``the drone actually did X'' is real and I have been careful not to over-claim it: my own wk22 capability audit, published in \texttt{docs/NEEDS\_VS\_CAPABILITIES.md}, scored simulation 7.7/10 and real-hardware 2.1/10 -- published internally before I wrote this report, precisely so that my reflection would be falsifiable against my own numbers rather than vibes. Leveson's point about safety being a \textit{system property}, not a component property \cite{leveson2011}, was the framing I used to decide which gaps mattered: a geofence that works in simulation but not in air is not a geofence, it is a placebo.

\textit{Evaluating my own performance against a criterion.} \rtag{M16.d-own} The criterion I will evaluate myself against is the one the project brief imposes on me: ten deliverables (D0--D7 plus the two internal milestones), AHEP4 M5/M7/M16/M17 evidence, and 12 Requirements. Judgement: on the ten deliverables I hit nine on time (D3 slipped by two days after flight day 2 cancelled) and contributed authored content to all ten. On the 12 Requirements I closed 12/12 in simulation, 0/12 verified on airborne hardware (because no flight occurred) -- a gap I report rather than hide. Cause: the flight cancellations were external, but the R01--R12 gap-close happened in a single wk22 session because I had deferred it; I chose not to do it earlier because I believed (wrongly) that hardware verification was imminent and would make bench-verification redundant. Lesson, applied this session: I rebuilt the test ladder so every R-requirement has a bench-executable test script by default, with the flight check as an upgrade, not a prerequisite. The falsifiable signal I will watch for in future projects is whether my compliance artefacts work \textit{before} hardware arrives, not after.

\textbf{3b. What I would add or refocus on.} \rtag{M16.c}\rtag{M17.d} My contribution pattern was dominated by \textit{building things}. I should shift it towards \textit{enabling people}, by roughly a third. Concretely:

\begin{itemize}
\item \textit{Pair-programming, not handoff.} My integration with Robin worked because of a clean interface, but Robin did not learn any more about vision from our collaboration. If I ran this again I would book two paired sessions per week where Robin and I worked on the same screen on vision code. The measurable test: after four weeks of pairing, could Robin independently diagnose the TFLite bbox coordinate bug I hit on 30~March? Today the answer is probably not. Next project, I want that to be yes.
\item \textit{Teaching-first documentation.} I wrote the \texttt{VISION\_PIPELINE\_FOR\_COLLEAGUE.md} and \texttt{COLLEAGUE\_CHECKLIST.md} in wk22, explicitly to onboard a next cohort. That was the right content, at the wrong time. The right time would have been wk14, with Robin or Edward as the test reader. I now see this is because teaching is a \textit{scheduling decision}, not a literary one: the quality of the prose matters much less than when it lands in the recipient's head.
\item \textit{Leading less by doing, more by proposing.} When Robin and I disagreed on ROS~2 vs raw MAVLink , I ``resolved'' it by writing a 20-line working heartbeat demo overnight. That was persuasion-by-fait-accompli, not leadership. A better version would have been a 15-minute whiteboard showing trade-offs \textit{before} writing any code. The transition from solo-problem-solver to team-enabler is the single skill schooling under-trains, and I can now name the pattern in myself. Naming the pattern is a prerequisite for changing it.
\end{itemize}

\textit{The tension between two values I did not previously notice.} I hold a commitment to \textit{delivering} and a commitment to \textit{developing people}. On this project, under time pressure, the first one consistently ate the second one. I had not previously noticed that these two values could conflict in me -- I thought of them as the same thing with different names. They are not. Recognising the tension is the reflection I take forward; resolving it is the self-development programme in \S4. The through-line: my technical output was high and my interpersonal leverage was low. In a project where the assessment rewards both, that is a real failure of allocation, and it is the main thing I take forward.
